\section{Related Work}
\label{sec:related}

\paragraph{Memory systems for LLM agents.}
Retrieval-augmented generation retrieves passages by embedding similarity
to the query \citep{lewis2020rag}, and most agent memory frameworks
inherit this ranking. MemGPT \citep{packer2023memgpt} treats context as a
paged virtual memory with explicit eviction, but eviction is driven by
recency and capacity pressure rather than a learned value. Generative
Agents \citep{park2023generative} introduce a retrieval score that blends
recency, importance, and relevance --- importantly, a \emph{multi-factor}
score --- but the importance term is produced ad hoc by an LLM rating
$1$--$10$, the weights are fixed and hand-set, and the score governs only
retrieval, not encoding depth or forgetting. MemoryBank
\citep{zhong2024memorybank} adds an Ebbinghaus-inspired decay so older
memories fade unless reinforced, again with hand-set dynamics. MemOS
\citep{memos2025} frames memory as an operating-system resource with
scheduling and lifecycle management, motivating a principled value signal
but not learning one. Across this line of work, the triage signal is
typically (i) single-factor (similarity or recency), or (ii)
multi-factor but hand-weighted and confined to retrieval. We differ on
both counts: our value is multi-factor, its weights are \emph{learned}
from a downstream objective, and the \emph{same} scalar drives encoding,
forgetting, and retrieval.

\paragraph{Value- and need-based retention in cognition.}
The factors we combine are not arbitrary; each corresponds to a robust
determinant of human retention. Value-directed remembering shows people
strategically retain high-value items and shed low-value detail
\citep{castel2008value}. Levels-of-processing
\citep{craik1972lop} and the self-reference effect \citep{rogers1977sre}
establish that semantically, goal-, and self-relevant encoding is more
durable than shallow encoding. Emotional arousal modulates consolidation
\citep{mcgaugh2000consolidation}. The rational analysis of memory
\citep{anderson1991adaptive} reframes forgetting itself as adaptive: the
retention function mirrors the environmental \emph{need-probability} of
information, and directed forgetting is a feature, not a bug
\citep{bjork1994memory}. Our contribution is to operationalise this
multi-determinant view as a single learned scalar for an artificial
agent, rather than to model human data --- a companion paper treats the
cognitive-modeling angle; here the target is engineering utility.

\paragraph{Learning what to remember.}
Reinforcement learning provides the natural framing for fitting a memory
policy to downstream outcome: the value of retaining an item is its
marginal contribution to future return \citep{sutton2018rl}. Because the
encode--forget--retrieve--answer pipeline is non-differentiable
(discrete keep/drop decisions, an external answerer), we fit weights with
a gradient-free black-box optimiser \citep{hansen2016cmaes} rather than
backpropagation. This keeps the value function low-dimensional (seven
weights) and interpretable, in contrast to end-to-end learned memory
controllers whose policies are opaque.

\paragraph{Evaluating memory.}
LongMemEval \citep{wu2025longmemeval} is the most directly relevant
benchmark: $500$ questions over long multi-session chat histories, with
gold-evidence turns flagged inside a haystack of distractor sessions. We
use it as our real-data testbed. Our methodological observation ---
that scoring relevance against the held-out question conflates retrieval
with retention, and that a forgetting policy must be evaluated
\emph{blind} to the future query --- applies to any retention metric
built on query-defined gold, and to our knowledge has not been made
explicit in this setting.
